% =============================================================================
% RESULTS
% Status: Updated for drone + ground camera analysis v0.4
% =============================================================================

% ============================================================================
\subsection{Part I: Drone Imagery Analysis}
% ============================================================================

\subsection{Drone Site 1: Santa Ines}

\subsubsection{Combined Spatial and Visual Analysis}

Figure~\ref{fig:si_combined} presents an integrated view of the Santa Ines drone survey analysis, demonstrating the fundamental difference between physical GPS locations and visual semantic clustering.

\begin{figure}[H]
    \centering
    \includegraphics[width=\textwidth]{si_combined_analysis.png}
    \caption{Combined analysis of Santa Ines drone imagery (97 images). \textbf{(A)} GPS coordinates overlaid on satellite imagery showing flight coverage; \textbf{(B)} UMAP projection of CLIP embeddings revealing two semantic clusters; \textbf{(C)} Representative close-up images capturing detailed crop inspection views; \textbf{(D)} Representative high-altitude overview images showing field row patterns.}
    \label{fig:si_combined}
\end{figure}

Panel A reveals the physical flight pattern over the Santa Ines field. The drone followed systematic transects across the crop rows, with close-up images (green markers) distributed along parallel flight lines and high-altitude overview images (red markers) captured at strategic positions. Geographically, both image types are \textbf{spatially intermixed}.

In contrast, Panel B shows the CLIP embedding space where images cluster by \textbf{visual content rather than geographic location}. Despite being captured from scattered GPS positions, all high-altitude images form a tight cluster clearly separated from the close-up images.

\subsubsection{Cluster Characteristics}

\begin{table}[H]
\centering
\caption{Santa Ines cluster statistics}
\label{tab:si_clusters}
\begin{tabular}{lcccc}
\toprule
\textbf{Category} & \textbf{Count} & \textbf{\%} & \textbf{UMAP X} & \textbf{UMAP Y} \\
\midrule
Close-up (crop inspection) & 89 & 91.8\% & $8.97 \pm 1.15$ & $2.15 \pm 1.28$ \\
High-altitude (overview) & 8 & 8.2\% & $11.30 \pm 0.70$ & $-0.56 \pm 0.27$ \\
\bottomrule
\end{tabular}
\end{table}

\begin{itemize}
    \item \textbf{Inter-cluster distance}: 3.56 UMAP units
    \item \textbf{Cross-category edges}: 8/291 (2.7\%)
    \item \textbf{Cluster cohesion}: Close-up 0.424, High-altitude 0.276
\end{itemize}

\subsubsection{Similarity Network}

\begin{figure}[H]
    \centering
    \includegraphics[width=0.9\textwidth]{si_similarity_network.png}
    \caption{Santa Ines similarity network. Gray edges connect images within the same category; yellow edges indicate cross-category connections.}
    \label{fig:si_network}
\end{figure}

% ============================================================================
\subsection{Drone Site 2: La Capilla}

\subsubsection{Combined Spatial and Visual Analysis}

Figure~\ref{fig:lc_combined} presents the La Capilla drone survey analysis, revealing a different clustering pattern than Santa Ines.

\begin{figure}[H]
    \centering
    \includegraphics[width=\textwidth]{lc_combined_analysis.png}
    \caption{Combined analysis of La Capilla drone imagery (111 images). \textbf{(A)} GPS coordinates showing flight coverage; \textbf{(B)} UMAP projection revealing two semantic clusters with large separation; \textbf{(C)} Representative very close-up images showing plant-level detail; \textbf{(D)} Representative medium-altitude images showing row structure.}
    \label{fig:lc_combined}
\end{figure}

The La Capilla dataset shows a striking spatial-semantic decoupling: The ``very close-up'' cluster (purple) represents images taken from a concentrated GPS region in the northwest of the field, while the ``medium altitude'' cluster (green) spans a larger geographic area. However, in embedding space, these form two completely separated clusters with \textbf{zero cross-category connections}.

\subsubsection{Visual Characteristics}

\textbf{Very close-up images} (26 images, 23.4\%) are characterized by:
\begin{itemize}
    \item Dense foliage filling the entire frame
    \item Individual leaf textures visible
    \item No visible row structure or soil
    \item Plant-level detail for crop inspection
\end{itemize}

\textbf{Medium-altitude images} (85 images, 76.6\%) exhibit:
\begin{itemize}
    \item Clear row structure with soil paths visible
    \item Multiple crop rows in each frame
    \item Geometric patterns from above
    \item Field-scale coverage for mapping
\end{itemize}

\subsubsection{Cluster Characteristics}

\begin{table}[H]
\centering
\caption{La Capilla cluster statistics}
\label{tab:lc_clusters}
\begin{tabular}{lcccc}
\toprule
\textbf{Category} & \textbf{Count} & \textbf{\%} & \textbf{UMAP X} & \textbf{UMAP Y} \\
\midrule
Very close-up (plant-level) & 26 & 23.4\% & $-0.43 \pm 0.34$ & $2.45 \pm 0.48$ \\
Medium altitude (row-visible) & 85 & 76.6\% & $12.07 \pm 1.85$ & $5.74 \pm 1.63$ \\
\bottomrule
\end{tabular}
\end{table}

\begin{itemize}
    \item \textbf{Inter-cluster distance}: 12.93 UMAP units
    \item \textbf{Cross-category edges}: 0/333 (0.0\%)
    \item \textbf{Cluster cohesion}: Very close-up 0.410, Medium altitude 0.726
\end{itemize}

\finding{La Capilla shows perfect cluster separation with zero cross-category edges, indicating strong visual distinction between image types}

\subsubsection{Similarity Network}

\begin{figure}[H]
    \centering
    \includegraphics[width=0.9\textwidth]{lc_similarity_network.png}
    \caption{La Capilla similarity network. All edges are within-category (gray), with no cross-category connections.}
    \label{fig:lc_network}
\end{figure}

\subsubsection{GPS vs Embedding Space}

Figure~\ref{fig:lc_gps} demonstrates the key finding: images that are geographically concentrated (very close-up cluster in GPS space) form a distinct semantic cluster in embedding space.

\begin{figure}[H]
    \centering
    \includegraphics[width=\textwidth]{lc_gps_vs_embeddings.png}
    \caption{La Capilla: GPS coordinates (left) vs CLIP embedding space (right). Color indicates longitude; red circles mark embedding outliers. The very close-up images are geographically concentrated but form a distinct semantic cluster.}
    \label{fig:lc_gps}
\end{figure}

% ============================================================================
\subsection{Cross-Site Comparison}

\subsubsection{Summary Statistics}

\begin{table}[H]
\centering
\caption{Cross-site comparison of clustering metrics}
\label{tab:comparison}
\begin{tabular}{lcc}
\toprule
\textbf{Metric} & \textbf{Santa Ines} & \textbf{La Capilla} \\
\midrule
Total images & 97 & 111 \\
Cluster ratio (primary/secondary) & 92\% / 8\% & 77\% / 23\% \\
Inter-cluster distance (UMAP units) & 3.56 & \textbf{12.93} \\
Cross-category edges (\%) & 2.7\% & \textbf{0.0\%} \\
Primary cluster cohesion & 0.424 & 0.726 \\
Secondary cluster cohesion & 0.276 & 0.410 \\
Semantic distinction & Close-up vs Overview & Plant-level vs Row-visible \\
\bottomrule
\end{tabular}
\end{table}

\subsubsection{Key Findings}

\begin{enumerate}
    \item \textbf{Consistent clustering}: Both sites showed clear bimodal clustering based on imaging altitude/purpose, despite different cluster ratios.

    \item \textbf{Stronger separation at La Capilla}: Inter-cluster distance was 3.6$\times$ larger at La Capilla (12.93 vs 3.56 units), reflecting a more pronounced visual distinction between image types.

    \item \textbf{Perfect separation achievable}: La Capilla achieved 0\% cross-category connections, demonstrating that embedding-based clustering can achieve complete category separation.

    \item \textbf{Cohesion patterns}: In both sites, the smaller cluster showed higher cohesion (tighter grouping), suggesting more homogeneous visual content in minority categories.
\end{enumerate}

\finding{CLIP embeddings consistently separate drone images by visual content (altitude/detail level) across different agricultural sites, without any domain-specific training}

\subsubsection{Cluster Cohesion Comparison}

\begin{figure}[H]
    \centering
    \begin{subfigure}[b]{0.48\textwidth}
        \includegraphics[width=\textwidth]{si_cluster_cohesion.png}
        \caption{Santa Ines}
    \end{subfigure}
    \hfill
    \begin{subfigure}[b]{0.48\textwidth}
        \includegraphics[width=\textwidth]{lc_cluster_cohesion.png}
        \caption{La Capilla}
    \end{subfigure}
    \caption{Cluster cohesion (mean distance to 5 nearest neighbors) at both sites.}
    \label{fig:cohesion_comparison}
\end{figure}

\observation{The secondary cluster (minority category) shows higher cohesion at both sites, suggesting that outlier image types have more consistent visual characteristics}

% ============================================================================
\subsection{Part II: Ground Camera Imagery Analysis}
% ============================================================================

Ground-level imagery captured with mobile phones (``Celular'') provides a complementary perspective to drone imagery, focusing on plant-level detail rather than field-scale patterns.

\subsubsection{Dataset Summary}

\begin{table}[H]
\centering
\caption{Ground camera datasets processed with CLIP embeddings}
\label{tab:ground_summary}
\begin{tabular}{lccccc}
\toprule
\textbf{Site} & \textbf{Images} & \textbf{Format} & \textbf{With GPS} & \textbf{Device} \\
\midrule
Santa Ines (Ground) & 1,325 & JPEG+HEIC & 1,325 (100\%) & Android \\
La Capilla (Ground) & 822 & JPEG & 799 (97\%) & Android \\
Apalta & 961 & HEIC & 960 (99.9\%) & iPhone XR \\
Camarico & 33 & JPEG & 27 (82\%) & Android \\
Lo Castillo & 1,521 & JPEG+HEIC & 1,521 (100\%) & Mixed \\
\midrule
\textbf{Total} & \textbf{4,662} & & \textbf{4,632 (99.4\%)} & \\
\bottomrule
\end{tabular}
\end{table}

\subsubsection{Embedding Space Analysis}

K-means clustering was applied to the UMAP-projected embeddings to identify natural groupings in the ground camera imagery:

\begin{table}[H]
\centering
\caption{Ground camera cluster analysis metrics}
\label{tab:ground_clusters}
\begin{tabular}{lcccc}
\toprule
\textbf{Site} & \textbf{Images} & \textbf{Optimal k} & \textbf{Silhouette} & \textbf{Inter-cluster Dist} \\
\midrule
Santa Ines (Ground) & 1,325 & 2 & 0.416 & 4.82 \\
La Capilla (Ground) & 822 & 2 & 0.472 & 3.50 \\
Apalta & 961 & 4 & 0.415 & 3.27 \\
Camarico & 33 & 2 & 0.461 & 2.35 \\
Lo Castillo & 1,521 & 2 & \textbf{0.537} & 5.38 \\
\bottomrule
\end{tabular}
\end{table}

\finding{Lo Castillo shows the strongest clustering among ground datasets (silhouette=0.537), followed by La Capilla (0.472). Lo Castillo also has the lowest cross-cluster rate (0.1\%) despite combining JPEG and HEIC formats}

% ----------------------------------------------------------------------------
\subsubsection{Ground Site 1: Santa Ines (Celular)}

Santa Ines ground camera is the largest ground dataset with 1,325 images combining both JPEG (635) and HEIC (690) formats. Figure~\ref{fig:sic_combined} presents the combined analysis.

\begin{figure}[H]
    \centering
    \includegraphics[width=\textwidth]{sic_combined_analysis.png}
    \caption{Combined analysis of Santa Ines ground camera imagery (1,325 images: 635 JPEG + 690 HEIC). \textbf{(A)} GPS coordinates showing field coverage; \textbf{(B)} UMAP projection revealing two clusters; \textbf{(C)} Cluster distribution by file format; \textbf{(D)} Dataset metrics showing 0.8\% cross-cluster edges.}
    \label{fig:sic_combined}
\end{figure}

\begin{figure}[H]
    \centering
    \begin{subfigure}[b]{0.48\textwidth}
        \includegraphics[width=\textwidth]{sic_similarity_network.png}
        \caption{Similarity network (0\% cross-cluster)}
    \end{subfigure}
    \hfill
    \begin{subfigure}[b]{0.48\textwidth}
        \includegraphics[width=\textwidth]{sic_temporal_analysis.png}
        \caption{Temporal clustering by survey date}
    \end{subfigure}
    \caption{Santa Ines ground camera: (a) k-NN similarity network showing perfect cluster separation with no cross-cluster edges; (b) Temporal analysis revealing clusters correspond to different survey dates (Jan 6 vs Jan 8, 2026).}
    \label{fig:sic_analysis}
\end{figure}

% ----------------------------------------------------------------------------
\subsubsection{Ground Site 2: La Capilla (Celular)}

\begin{figure}[H]
    \centering
    \includegraphics[width=\textwidth]{lcc_combined_analysis.png}
    \caption{Combined analysis of La Capilla ground camera imagery (822 images). Unlike the drone dataset which showed perfect separation, ground camera clustering shows moderate overlap with 1.5\% cross-cluster edges.}
    \label{fig:lcc_combined}
\end{figure}

\begin{figure}[H]
    \centering
    \begin{subfigure}[b]{0.48\textwidth}
        \includegraphics[width=\textwidth]{lcc_similarity_network.png}
        \caption{Similarity network (1.5\% cross-cluster)}
    \end{subfigure}
    \hfill
    \begin{subfigure}[b]{0.48\textwidth}
        \includegraphics[width=\textwidth]{lcc_cluster_cohesion.png}
        \caption{Cluster cohesion comparison}
    \end{subfigure}
    \caption{La Capilla ground camera analysis showing moderate cluster separation.}
    \label{fig:lcc_analysis}
\end{figure}

% ----------------------------------------------------------------------------
\subsubsection{Ground Site 3: Apalta}

Apalta represents the largest ground camera dataset (961 images) and shows more complex clustering with 4 distinct groups.

\begin{figure}[H]
    \centering
    \includegraphics[width=\textwidth]{ap_combined_analysis.png}
    \caption{Combined analysis of Apalta ground camera imagery (961 images). K-means identified 4 clusters, suggesting greater visual diversity in this dataset.}
    \label{fig:ap_combined}
\end{figure}

\begin{figure}[H]
    \centering
    \begin{subfigure}[b]{0.48\textwidth}
        \includegraphics[width=\textwidth]{ap_similarity_network.png}
        \caption{Similarity network (2.0\% cross-cluster)}
    \end{subfigure}
    \hfill
    \begin{subfigure}[b]{0.48\textwidth}
        \includegraphics[width=\textwidth]{ap_cluster_cohesion.png}
        \caption{Cluster cohesion (4 clusters)}
    \end{subfigure}
    \caption{Apalta ground camera analysis showing 4-cluster structure with varying cohesion.}
    \label{fig:ap_analysis}
\end{figure}

\observation{Apalta's 4-cluster structure (vs. 2 clusters at other sites) suggests greater internal visual diversity, possibly due to varied imaging conditions or crop heterogeneity across the field}

% ----------------------------------------------------------------------------
\subsubsection{Ground Site 4: Camarico}

Camarico represents the smallest dataset (33 images), providing a test of methodology robustness on limited data.

\begin{figure}[H]
    \centering
    \includegraphics[width=0.8\textwidth]{cam_combined_analysis.png}
    \caption{Combined analysis of Camarico ground camera imagery (33 images). Despite the small sample size, meaningful clustering is observed with silhouette=0.461.}
    \label{fig:cam_combined}
\end{figure}

\finding{Even with only 33 images, Camarico produced interpretable clustering (silhouette=0.461, 6.1\% cross-cluster edges), demonstrating methodology applicability to minimal datasets}

% ----------------------------------------------------------------------------
\subsubsection{Ground Site 5: Lo Castillo}

Lo Castillo represents the largest ground camera dataset (1,521 images), combining both JPEG (1,215) and HEIC (306) formats. This site provides a comprehensive test of the methodology on mixed-format, large-scale ground imagery.

\begin{figure}[H]
    \centering
    \includegraphics[width=\textwidth]{loc_combined_analysis.png}
    \caption{Combined analysis of Lo Castillo ground camera imagery (1,521 images: 1,215 JPEG + 306 HEIC). \textbf{(A)} GPS coordinates showing field coverage with 100\% GPS extraction success; \textbf{(B)} UMAP projection revealing two clusters with strong separation; \textbf{(C)} Format distribution by cluster showing both JPEG and HEIC images present in each cluster; \textbf{(D)} Dataset metrics showing silhouette=0.537 and only 0.1\% cross-cluster edges.}
    \label{fig:loc_combined}
\end{figure}

\begin{figure}[H]
    \centering
    \begin{subfigure}[b]{0.48\textwidth}
        \includegraphics[width=\textwidth]{loc_similarity_network.png}
        \caption{Similarity network (0.1\% cross-cluster)}
    \end{subfigure}
    \hfill
    \begin{subfigure}[b]{0.48\textwidth}
        \includegraphics[width=\textwidth]{loc_cluster_cohesion.png}
        \caption{Cluster cohesion comparison}
    \end{subfigure}
    \caption{Lo Castillo ground camera analysis: (a) k-NN similarity network showing near-perfect cluster separation with only 4 cross-cluster edges out of 4,563 total; (b) Cluster cohesion showing Cluster 1 is tighter despite having fewer samples.}
    \label{fig:loc_analysis}
\end{figure}

\subsubsection{Lo Castillo Cluster Characteristics}

\begin{table}[H]
\centering
\caption{Lo Castillo cluster statistics}
\label{tab:loc_clusters}
\begin{tabular}{lcccc}
\toprule
\textbf{Category} & \textbf{Count} & \textbf{\%} & \textbf{UMAP X} & \textbf{UMAP Y} \\
\midrule
Cluster 0 (Primary) & 1,186 & 78.0\% & $7.47 \pm 1.42$ & $7.14 \pm 1.68$ \\
Cluster 1 (Secondary) & 335 & 22.0\% & $2.26 \pm 1.00$ & $5.79 \pm 0.81$ \\
\bottomrule
\end{tabular}
\end{table}

\begin{itemize}
    \item \textbf{Inter-cluster distance}: 5.38 UMAP units
    \item \textbf{Cross-cluster edges}: 4/4,563 (\textbf{0.1\%})---lowest among all ground camera datasets
    \item \textbf{Silhouette score}: 0.537---\textbf{highest among all ground camera datasets}
    \item \textbf{GPS coverage}: 100\% (all 1,521 images have valid GPS coordinates)
\end{itemize}

\finding{Lo Castillo achieves the best clustering performance among all ground camera datasets, with the highest silhouette score (0.537) and lowest cross-cluster rate (0.1\%), despite being the largest dataset (1,521 images)}

\subsubsection{Multi-Format Analysis}

Unlike other ground camera datasets which used a single file format, Lo Castillo combines images from multiple devices:

\begin{table}[H]
\centering
\caption{Lo Castillo format distribution by cluster}
\label{tab:loc_format}
\begin{tabular}{lccc}
\toprule
\textbf{Format} & \textbf{Total} & \textbf{Cluster 0} & \textbf{Cluster 1} \\
\midrule
JPEG & 1,215 (79.9\%) & 952 (80.3\%) & 263 (78.5\%) \\
HEIC & 306 (20.1\%) & 234 (19.7\%) & 72 (21.5\%) \\
\bottomrule
\end{tabular}
\end{table}

\observation{Both file formats are distributed proportionally across clusters, indicating that CLIP embeddings cluster by visual content rather than file format or device characteristics}

% ----------------------------------------------------------------------------
\subsubsection{Ground Camera: Network Analysis Summary}

\begin{table}[H]
\centering
\caption{Ground camera similarity network analysis (k=3 nearest neighbors)}
\label{tab:ground_network}
\begin{tabular}{lcccc}
\toprule
\textbf{Site} & \textbf{Within-cluster} & \textbf{Cross-cluster} & \textbf{Cross \%} & \textbf{Separation} \\
\midrule
Santa Ines (Ground) & 3,942 & 33 & 0.8\% & Strong \\
La Capilla (Ground) & 2,428 & 38 & 1.5\% & Strong \\
Apalta & 2,824 & 59 & 2.0\% & Strong \\
Camarico & 93 & 6 & 6.1\% & Moderate \\
Lo Castillo & 4,559 & 4 & \textbf{0.1\%} & Very Strong \\
\bottomrule
\end{tabular}
\end{table}

\finding{All ground camera datasets achieved strong cluster separation ($<$2\% cross-cluster edges), with Lo Castillo achieving the lowest cross-cluster rate (0.1\%) and highest silhouette score (0.537) despite being the largest dataset (1,521 images)}

% ----------------------------------------------------------------------------
\subsubsection{Temporal Clustering Discovery}

Analysis of the JPEG subset (635 images with extractable timestamps) revealed temporal patterns in the Santa Ines ground camera data:

\begin{itemize}
    \item \textbf{January 8, 2026}: 526 images (83\%)
    \item \textbf{January 6, 2026}: 109 images (17\%)
\end{itemize}

The 690 HEIC images (iPhone format) show different naming conventions, but the combined JPEG+HEIC dataset maintains strong cluster separation (0.8\% cross-cluster edges).

\observation{CLIP embeddings detect visual differences that may correlate with survey dates, lighting conditions, or subtle environmental changes---useful for quality control in multi-day surveys}

This temporal sensitivity demonstrates that vision-language embeddings can serve as an implicit indicator of survey conditions, potentially useful for:
\begin{itemize}
    \item Quality control (detecting inconsistent imaging conditions)
    \item Change detection (monitoring crop development over time)
    \item Data stratification (grouping images by visual similarity for analysis)
\end{itemize}

\subsubsection{Comparison: Drone vs Ground Camera}

\begin{table}[H]
\centering
\caption{Clustering comparison across capture modalities}
\label{tab:modality_comparison}
\begin{tabular}{llccc}
\toprule
\textbf{Site} & \textbf{Modality} & \textbf{Images} & \textbf{Silhouette} & \textbf{Inter-dist} \\
\midrule
Santa Ines & Drone & 97 & --- & 3.56 \\
Santa Ines & Ground & 1,325 & 0.416 & 4.82 \\
La Capilla & Drone & 111 & --- & 12.93 \\
La Capilla & Ground & 822 & 0.472 & 3.50 \\
\bottomrule
\end{tabular}
\end{table}

Santa Ines and La Capilla ground datasets show similar inter-cluster distances (4.82 and 3.50), while their drone datasets show different patterns (3.56 vs 12.93). This suggests that ground camera datasets capture more uniform visual content (plant-level views), while drone datasets vary based on flight altitude patterns.

\subsubsection{Key Characteristics}

Ground camera imagery differs fundamentally from drone imagery:

\begin{itemize}
    \item \textbf{Perspective}: Horizontal view at plant level vs. top-down aerial view
    \item \textbf{Content}: Individual plants, leaves, fruits, and row structure visible from ground level
    \item \textbf{Scale}: Much larger datasets (33--1,521 images per site vs. 97--111 for drone)
    \item \textbf{Collection method}: Walking through fields vs. programmed flight paths
    \item \textbf{Multi-format support}: Successfully processed both JPEG (Android) and HEIC (iPhone) formats
\end{itemize}

\finding{Ground camera datasets contain up to 15$\times$ more images than drone datasets (e.g., Lo Castillo with 1,521 images), reflecting the intensive ground-truthing effort typical in precision agriculture research}

\subsubsection{GPS Coverage}

GPS extraction was successful for the vast majority of ground camera images:

\begin{itemize}
    \item \textbf{Android phones}: Standard EXIF GPS extraction (97--100\% success)
    \item \textbf{iPhone (HEIC)}: Required specialized IFD-based GPS extraction (99.9\% success)
\end{itemize}

\subsubsection{Geospatial Outputs}

For each ground camera dataset, the following files were generated:

\begin{itemize}
    \item \textbf{KML/KMZ}: Photo locations viewable in Google Earth
    \item \textbf{Shapefile (UTM 19S)}: GIS-compatible format with UTM coordinates
    \item \textbf{CSV}: Complete dataset with UMAP embeddings, GPS, and metadata
\end{itemize}

% ============================================================================
\subsection{Part III: Combined Dataset Summary}
% ============================================================================

\begin{table}[H]
\centering
\caption{Complete dataset summary across all capture modalities}
\label{tab:all_datasets}
\begin{tabular}{llccc}
\toprule
\textbf{Site} & \textbf{Modality} & \textbf{Images} & \textbf{GPS Coverage} & \textbf{FiftyOne Dataset} \\
\midrule
Santa Ines & Drone & 97 & 100\% & santa\_ines\_dron \\
Santa Ines & Ground & 1,325 & 100\% & santa\_ines\_celular \\
La Capilla & Drone & 111 & 100\% & la\_capilla\_dron \\
La Capilla & Ground & 822 & 97\% & la\_capilla\_celular \\
Apalta & Ground & 961 & 99.9\% & apalta \\
Camarico & Ground & 33 & 82\% & camarico \\
Lo Castillo & Ground & 1,521 & 100\% & lo\_castillo \\
\midrule
\multicolumn{2}{l}{\textbf{Total}} & \textbf{4,870} & \textbf{99.4\%} & \\
\bottomrule
\end{tabular}
\end{table}

\finding{The combined multi-modal dataset comprises 4,870 agricultural images across 5 sites and 2 capture modalities, all processed with identical CLIP embedding methodology}

\subsubsection{Data Availability by Site and Modality}

An important observation is that not all sites have both capture modalities available:

\begin{table}[H]
\centering
\caption{Modality availability comparison across sites}
\label{tab:modality_availability}
\begin{tabular}{lccccc}
\toprule
\textbf{Site} & \textbf{Drone} & \textbf{Ground} & \textbf{Drone Images} & \textbf{Ground Images} & \textbf{Ratio (G:D)} \\
\midrule
Santa Ines & \checkmark & \checkmark & 97 & 1,325 & 13.7:1 \\
La Capilla & \checkmark & \checkmark & 111 & 822 & 7.4:1 \\
Apalta & --- & \checkmark & 0 & 961 & Ground only \\
Camarico & --- & \checkmark & 0 & 33 & Ground only \\
Lo Castillo & --- & \checkmark & 0 & 1,521 & Ground only \\
\bottomrule
\end{tabular}
\end{table}

\observation{Two sites (Santa Ines and La Capilla) have complete multi-modal coverage with both drone and ground camera imagery, while three sites (Apalta, Camarico, and Lo Castillo) have only ground camera data available}

Key findings on data availability:

\begin{itemize}
    \item \textbf{Multi-modal sites}: Santa Ines and La Capilla have both drone and ground camera imagery, enabling direct comparison of embedding behavior across capture modalities
    \item \textbf{Ground-only sites}: Apalta, Camarico, and Lo Castillo were surveyed only with ground cameras, limiting analysis to single-modality embedding patterns
    \item \textbf{Ground camera ratio}: For sites with both modalities, ground camera images outnumber drone images by 7--14$\times$
    \item \textbf{Largest dataset}: Lo Castillo has the largest ground camera dataset (1,521 images), with the best clustering performance (silhouette=0.537)
    \item \textbf{Smallest dataset}: Camarico has the smallest dataset (33 images), significantly smaller than other ground camera datasets (822--1,521 images)
\end{itemize}

\finding{Camarico represents a minimal-data scenario with only 33 ground camera images and no drone coverage---useful for testing methodology robustness on limited data but insufficient for detailed cluster analysis}

\subsubsection{Drone vs Ground Camera: Side-by-Side Comparison}

For the two sites with complete multi-modal coverage, we can directly compare embedding behavior:

\begin{table}[H]
\centering
\caption{Detailed comparison of drone vs ground camera for multi-modal sites}
\label{tab:drone_vs_ground}
\begin{tabular}{llcccc}
\toprule
\textbf{Site} & \textbf{Modality} & \textbf{Images} & \textbf{Clusters} & \textbf{Inter-dist} & \textbf{Primary Driver} \\
\midrule
Santa Ines & Drone & 97 & 2 & 3.56 & Altitude \\
Santa Ines & Ground & 1,325 & 2 & 4.82 & Mixed format \\
\midrule
La Capilla & Drone & 111 & 2 & 12.93 & Altitude \\
La Capilla & Ground & 822 & 2 & 3.50 & Visual content \\
\bottomrule
\end{tabular}
\end{table}

\observation{Ground camera datasets show consistent inter-cluster distances (3.50--4.82), while drone datasets vary more widely (3.56--12.93). This suggests ground-level imagery captures more uniform visual content regardless of site, while aerial imagery varies based on flight patterns and altitude choices}
